% !TeX program = lualatex
% ================================================================
% Chinese Mathematical Dictionary
%
% Generated by the server using the following settings:
%
%   language            Chinese (zho)
%   text direction      left to right
%   font                Noto Sans CJK SC
%   fallback font       Noto Serif
%   built from          latex/dictionary/mathematicalDictionary.tex
%   tier 3 vocabulary   green   RGB 0 128 0
%   English text        black   RGB 0 0 0
%   Chinese text        purple  RGB 112 48 160
%   layout macros       version 3
%
% Please don't edit any information in this comment block - the server
% compares it against the current settings and rebuilds this file when
% they differ, so changes here would be overwritten. However, feel free
% to make updates to the LaTeX below if you want to edit it.
% ================================================================
%
% ---- What is safe to edit in this file ----------------------------------
%
% Safe:     the Chinese word and the Chinese definition - the third
%           and fourth arguments of each entry. That is what this file is for:
%           if a translation reads wrongly to somebody who speaks the language,
%           correct it here and every sheet that uses the word follows.
%
% Not safe: the English word and the English definition - the first and second
%           arguments. The first is how an entry is matched to the shared
%           dictionary; the second is the wording the translation was made
%           from, and is how the server knows the translation is still current.
%           Changing an English definition here does not change the shared
%           dictionary - it only makes this entry look up to date when it is
%           not. Edit the English in the shared dictionary instead, and this
%           file will be brought back into step on the next run.
%
% Entries are added and removed by the server, following the shared dictionary.
% An entry the server cannot read is reported rather than guessed at, so if a
% run complains about this file, look for a missing brace.
% -------------------------------------------------------------------------

\documentclass[11pt]{article}
\usepackage[a4paper,margin=18mm]{geometry}
% Characters this language's font has not got are borrowed from another,
% rather than being silently dropped - see docs/translated-sheet-latex.md
\directlua{%
  if luaotfload and luaotfload.add_fallback then
    luaotfload.add_fallback("eallatinfallback",
      { "Noto Serif:mode=harf;" })
  else
    tex.error("this luaotfload is too old to register a fallback font")
  end
}
\usepackage[english]{babel}
\babelprovide[import]{chinese}
\babelfont[chinese]{rm}[Renderer=Harfbuzz, BoldFont={Noto Sans CJK SC}, ItalicFont={Noto Sans CJK SC}, BoldItalicFont={Noto Sans CJK SC}, RawFeature={fallback=eallatinfallback}]{Noto Sans CJK SC}
\babelfont[chinese]{sf}[Renderer=Harfbuzz, BoldFont={Noto Sans CJK SC}, ItalicFont={Noto Sans CJK SC}, BoldItalicFont={Noto Sans CJK SC}, RawFeature={fallback=eallatinfallback}]{Noto Sans CJK SC}
\newcommand{\ealtext}[1]{\foreignlanguage{chinese}{#1}}
\newcommand{\ealtextblock}[1]{{\raggedright\foreignlanguage{chinese}{#1}\par}}
% ================================================================
% EAL layout helpers
% ================================================================
\usepackage{xcolor}

\definecolor{ealkeycolour}{RGB}{0,128,0}
\definecolor{ealenglish}{RGB}{0,0,0}
\definecolor{ealtwocolour}{RGB}{112,48,160}

\newsavebox{\ealboxen}
\newsavebox{\ealboxtr}
\newlength{\ealglosswd}
\newlength{\ealglossraise}
\setlength{\ealglossraise}{1.15em}

% a tier 3 word, in English and in translation alike
\newcommand{\ealkey}[1]{\textcolor{ealkeycolour}{#1}}
\newcommand{\ealkeytr}[1]{\textcolor{ealkeycolour}{\ealtext{#1}}}

% One translated word above one English word. Built from kernel
% primitives, never a tabular, which would break wherever the sheet
% loads array, booktabs or tabularx. The box is as wide as the wider of
% the two, so a long translation pushes its neighbours aside instead of
% overlapping them, and it claims its full height so a table row or a
% TikZ node makes room for it.
\newcommand{\ealgloss}[2]{%
  \sbox{\ealboxen}{\ealkey{#1}}%
  \sbox{\ealboxtr}{\small\textcolor{ealtwocolour}{\ealtext{#2}}}%
  \setlength{\ealglosswd}{\wd\ealboxen}%
  \ifdim\wd\ealboxtr>\ealglosswd\setlength{\ealglosswd}{\wd\ealboxtr}\fi
  \leavevmode
  \makebox[\ealglosswd][c]{%
    \makebox[0pt][c]{\usebox{\ealboxen}}%
    \makebox[0pt][c]{\raisebox{\ealglossraise}{\usebox{\ealboxtr}}}%
  }%
}

% opens up line spacing around a block of glosses, so the raised
% translations sit clear of the line above
\newenvironment{ealglossed}{\par\linespread{2.1}\selectfont\ignorespaces}{\par}

% a whole translated sentence above its English counterpart. block level,
% so it does not belong inside a tabular cell
\newcommand{\ealpara}[2]{%
  \par\addvspace{0.5\baselineskip}%
  {\color{ealtwocolour}\ealtextblock{#1}}%
  \penalty200
  {\color{ealenglish}#2\par}%
  \addvspace{0.5\baselineskip}%
}

% ================================================================
% Dictionary layout
% ================================================================
\usepackage{multicol}

\newlength{\ealdictgap}
\setlength{\ealdictgap}{1.2mm}

% english word / english meaning / translated word / translated meaning
\newcommand{\dictentrytr}[4]{%
  \par\addvspace{\ealdictgap}%
  \noindent
  {\bfseries\ealkey{#1}}\quad{\small #2}\par
  \nopagebreak
  {\color{ealtwocolour}\ealtext{{\bfseries #3}~\textemdash{} {\small #4}}\par}%
  \par
}

\pagestyle{empty}
\setlength{\parindent}{0pt}

\begin{document}
{\large\bfseries Chinese Mathematical Dictionary}\par\medskip

\dictentrytr{adjacent}{next to, and sharing a side or a corner}{相邻的}{挨着，并共有一条边或一个角}
\dictentrytr{arc}{part of the curved edge of a circle}{弧}{圆的弯曲边缘的一部分}
\dictentrytr{area}{the amount of space inside a flat shape}{面积}{平面图形内部空间的大小}
\dictentrytr{average}{a single value chosen to stand for a whole set of data}{平均数}{被选中用来代表整组数据的单个值}
\dictentrytr{axis}{one of the lines a graph is measured against}{坐标轴}{图像用来作为测量基准的线之一}
\dictentrytr{bearing}{a direction given as an angle clockwise from north, written with three figures}{方位角}{从正北方向顺时针量出的角所表示的方向，用三位数字写出}
\dictentrytr{bisector}{a line that cuts an angle or another line exactly in half}{平分线}{把角或另一条线恰好分成两半的线}
\dictentrytr{chord}{a straight line joining two points on the edge of a circle}{弦}{连接圆边缘上两点的直线}
\dictentrytr{circumference}{the distance all the way round a circle}{圆周长}{绕圆一周的距离}
\dictentrytr{coefficient}{the number in front of a letter, showing how many of it there are}{系数}{字母前面的数，表示有多少个该字母}
\dictentrytr{congruent}{exactly the same shape and size}{全等的}{形状和大小完全相同}
\dictentrytr{correlation}{how closely two sets of data are related}{相关性}{两组数据相互关联的密切程度}
\dictentrytr{cumulative}{adding up as you go along, so each total includes the ones before}{累积的}{一边进行一边相加，因此每个总数都包括前面的数}
\dictentrytr{data}{the set of values collected}{数据}{收集到的一组数值}
\dictentrytr{decimal}{a number written with a point, where the digits after the point are parts of a whole}{小数}{用小数点写出的数，其中小数点后的数字表示整体的部分}
\dictentrytr{denominator}{the number below the line in a fraction, showing how many equal parts the whole is split into}{分母}{分数中线下面的数，表示整体被分成多少个相等的部分}
\dictentrytr{diameter}{a straight line across a circle through its centre}{直径}{穿过圆心、横跨圆的直线}
\dictentrytr{difference}{the answer when one number is subtracted from another}{差}{一个数减去另一个数所得的结果}
\dictentrytr{digit}{a single figure from 0 to 9 that a number is written with}{数字}{用来书写数的一个单个数字，从0到9}
\dictentrytr{edge}{a straight line where two faces of a solid meet}{棱}{立体图形中两个面相交的直线}
\dictentrytr{equivalent}{having the same value, even though it is written differently}{等值}{具有相同的值，即使写法不同}
\dictentrytr{estimate}{a sensible answer worked out quickly from rounded numbers}{估算}{利用取整后的数快速算出的合理答案}
\dictentrytr{expand}{to multiply out brackets so the expression has no brackets left}{展开}{把括号乘开，使表达式不再含括号}
\dictentrytr{expression}{a group of numbers and letters joined by operations, with no equals sign}{表达式}{由数和字母通过运算连接成的一组，没有等号}
\dictentrytr{face}{a flat surface of a solid shape}{面}{立体图形的一个平面}
\dictentrytr{factor}{a whole number that divides exactly into another number}{因数}{能整除另一个数的整数}
\dictentrytr{factorise}{to write an expression as a product, usually by putting brackets back in}{因式分解}{把一个表达式写成乘积形式，通常是把括号放回去}
\dictentrytr{formula}{a rule written with letters, showing how to work one quantity out from others}{公式}{用字母写成的规则，说明如何从其他量算出一个量}
\dictentrytr{fraction}{a number written as one whole split into equal parts, with a number of those parts taken}{分数}{写成一个整体分成相等的部分，并取其中若干部分的数}
\dictentrytr{frequency}{how many times something happens}{频数}{某事发生的次数}
\dictentrytr{gradient}{how steep a line is, found by how far it rises for each step across}{斜率}{一条线有多陡，通过每横向移动一步它上升多少来求得}
\dictentrytr{highest common factor}{the largest whole number that divides exactly into two or more numbers}{最大公因数}{能整除两个或更多数的最大整数}
\dictentrytr{hypotenuse}{the longest side of a right-angled triangle, opposite the right angle}{斜边}{直角三角形中最长的边，与直角相对}
\dictentrytr{identity}{a statement that is true for every value of the letters in it}{恒等式}{对其中字母的每一个值都成立的陈述}
\dictentrytr{index}{the small raised number showing what power something is raised to}{指数}{表示某数被提升到多少次幂的小的上标数字}
\dictentrytr{inequality}{a statement that one quantity is larger or smaller than another}{不等式}{表示一个量大于或小于另一个量的陈述}
\dictentrytr{integer}{a whole number, which may be positive, negative or zero}{整数}{一个整数，可以是正数、负数或零}
\dictentrytr{intercept}{the point where a graph crosses an axis}{截距}{图像与坐标轴相交的点}
\dictentrytr{linear}{making a straight line when drawn as a graph}{线性的}{画成图像时形成一条直线}
\dictentrytr{lowest common multiple}{the smallest number that two or more numbers all divide into exactly}{最小公倍数}{两个或更多数都能整除的最小数}
\dictentrytr{multiple}{a number you get by multiplying another number by a whole number}{倍数}{用一个整数乘另一个数所得的数}
\dictentrytr{multiply}{to add a number to itself a given number of times}{乘}{把一个数自身相加给定的次数}
\dictentrytr{nth term}{a rule that gives any number in a sequence from its position}{通项公式}{根据位置给出数列中任意一项的规则}
\dictentrytr{numerator}{the number above the line in a fraction, showing how many equal parts are taken}{分子}{分数中线上面的数，表示取了多少个相等的部分}
\dictentrytr{parallel}{always the same distance apart and never meeting}{平行的}{始终相距相同距离且永不相交}
\dictentrytr{percentage}{a number of parts out of a hundred}{百分数}{表示一百份中占多少份的数}
\dictentrytr{perimeter}{the total distance around the edge of a shape}{周长}{图形边缘一周的总距离}
\dictentrytr{perpendicular}{meeting at a right angle}{垂直的}{以直角相交}
\dictentrytr{place value}{the value a digit has because of where it sits in a number}{位值}{一个数字由于它在数中的位置而具有的值}
\dictentrytr{polygon}{a flat shape with straight sides}{多边形}{有直边的平面图形}
\dictentrytr{power}{how many times a number is multiplied by itself}{幂}{一个数被自身相乘的次数}
\dictentrytr{prime}{a whole number greater than 1 that divides only by itself and 1}{质数}{大于1且只能被自身和1整除的整数}
\dictentrytr{product}{the answer when two or more numbers are multiplied}{积}{两个或多个数相乘所得的结果}
\dictentrytr{proportion}{a relationship where two amounts change by the same factor}{比例}{两个量按相同倍数变化的关系}
\dictentrytr{quadratic}{an expression or equation whose highest power is a square}{二次的}{最高次幂为平方的表达式或方程}
\dictentrytr{quotient}{the answer when one number is divided by another}{商}{一个数除以另一个数所得的结果}
\dictentrytr{radius}{the distance from the centre of a circle to its edge}{半径}{从圆心到圆边缘的距离}
\dictentrytr{ratio}{a way of comparing two or more amounts, showing how much of one there is for each of the other}{比}{一种比较两个或更多量的方式，表示相对于另一个量，其中一个量有多少}
\dictentrytr{rearrange}{to change the order of a formula so a different letter is on its own}{变形}{改变公式的排列顺序，使另一个字母单独在一边}
\dictentrytr{reciprocal}{the number you multiply by to get 1, found by turning a fraction upside down}{倒数}{与你相乘得到1的数，可通过把分数上下颠倒求得}
\dictentrytr{recur}{to repeat the same digits forever after the decimal point}{循环}{小数点后永远重复相同的数字}
\dictentrytr{root}{a number that gives another number when multiplied by itself a given number of times}{根}{一个数，当它被自身相乘给定次数时得到另一个数}
\dictentrytr{round}{to replace a number with a nearby simpler one}{四舍五入}{用一个附近的更简单的数代替一个数}
\dictentrytr{scale factor}{the number every length is multiplied by when a shape is enlarged}{比例因子}{图形放大时每条长度都乘以的数}
\dictentrytr{sector}{a slice of a circle between two radii}{扇形}{圆中两条半径之间的一块}
\dictentrytr{segment}{the part of a circle cut off by a chord}{弓形}{圆被弦截下的部分}
\dictentrytr{sequence}{a list of numbers that follows a rule}{数列}{按规则排列的一列数}
\dictentrytr{similar}{the same shape but a different size}{相似的}{形状相同但大小不同}
\dictentrytr{simplify}{to write something in its smallest or plainest form without changing its value}{化简}{在不改变值的情况下，把某式写成最小或最简形式}
\dictentrytr{substitute}{to put a number in place of a letter}{代入}{用一个数代替一个字母}
\dictentrytr{sum}{the answer when numbers are added together}{和}{若干个数相加所得的结果}
\dictentrytr{surd}{a root that cannot be written exactly as a fraction, so it is left in root form}{不尽根}{不能精确地写成分数、因此保留根号形式的根}
\dictentrytr{tangent}{a line that touches a curve at exactly one point}{切线}{与曲线恰好只在一个点接触的直线}
\dictentrytr{term}{one part of an expression or one number in a sequence}{项}{表达式的一部分，或数列中的一个数}
\dictentrytr{transformation}{a change to a shape's position, size or orientation}{变换}{图形位置、大小或方向的变化}
\dictentrytr{unitary}{a method that first finds the value of one part or one item}{归一法}{先求出一份或一个物品的值的方法}
\dictentrytr{variable}{a letter standing for a quantity that can change}{变量}{代表可以变化的量的字母}
\dictentrytr{vector}{a quantity with both size and direction}{向量}{既有大小又有方向的量}
\dictentrytr{vertex}{a corner where two or more edges meet}{顶点}{两条或更多条棱相交处的角}
\dictentrytr{volume}{the amount of space inside a solid shape}{体积}{立体图形内部空间的大小}

\end{document}
